\section[Traitements simples de données textuelles]{Traitements simples de données textuelles}

\subsection[Introduction]{Introduction}

\begin{frame}{Contenu}
	Ce cours n'a pas vocation à être très avancé sur le traitement des données textuelles. Mais nous introduirons les concepts suivants:
	\begin{itemize}
		\item Tokenisation et mots vides (stop words)
		\item Lemmatisation et analyse grammaticale (Part of Speech)
		\item Vectorisation
		\item TF-IDF (term frequency, inverse document frequency)
	\end{itemize}
	Nous illustrerons ces éléments sur le dataset d'analyse de sentiment \href{https://www.kaggle.com/datasets/yasserh/imdb-movie-ratings-sentiment-analysis}{IMDb}. Il existe des librairies disposant de fonctions pour traiter les textes en langue anglaise, comme \href{https://www.nltk.org/}{NLTK (Natural Language Toolkit)}.
\end{frame}

\subsection[Tokenisation]{Tokenisation}

\begin{frame}[fragile]{Tokenisation}
	La tokenisation consiste à découper une phrase en éléments simples, \textbf{tokens}. On en profite en général pour passer le texte en lettres minuscules et pour supprimer la ponctuation.\\
	
	Par exemple:
	\begin{lstlisting}
		reviews = [
		"I loved every second of this masterpiece!",
		"The actors were terrible and the story was boring.",
		"Amazing music and great special effects!",
		"I would rather watch paint dry than see this movie.",
		"A complete waste of time and money!",
		"Don't miss this masterpiece!"
		]
	\end{lstlisting}
\end{frame}

\begin{frame}[fragile]{Tokenisation}
	Donnerait les tokens suivants:
	\begin{lstlisting}
		['i', 'loved', 'every', 'second', 'of', 'this', 'masterpiece']
		['the', 'actors', 'were', 'terrible', 'and', 'the', 'story', 'was', 'boring']
		['amazing', 'music', 'and', 'great', 'special', 'effects']
		['i', 'would', 'rather', 'watch', 'paint', 'dry', 'than', 'see', 'this', 'movie']
		['a', 'complete', 'waste', 'of', 'time', 'and', 'money']
		['do', "n't", 'miss', 'this', 'masterpiece']
	\end{lstlisting}
	\vfill
	Méthodes: \href{https://www.nltk.org/api/nltk.tokenize.html\#nltk-tokenize-package}{\texttt{nltk.word\_tokenize}}
\end{frame}

\begin{frame}[fragile]{Tokenisation - Stop Words}
	NLTK dispose d'une liste de mots vides (stop words) qui n'apportent que peu ou pas de sens à la phrase. Ou alors qui sont tellement communs qu'on les retrouverait à la fois dans des critiques positives et négatives:
	\begin{lstlisting}
		['loved', 'every', 'second', 'masterpiece']
		['actors', 'terrible', 'story', 'boring']
		['amazing', 'music', 'great', 'special', 'effects']
		['would', 'rather', 'watch', 'paint', 'dry', 'see', 'movie']
		['complete', 'waste', 'time', 'money']
		["n't", 'miss', 'masterpiece']
	\end{lstlisting}
	\vfill
	Méthodes: \href{https://www.nltk.org/howto/corpus.html\#word-lists-and-lexicons}{\texttt{nltk.corpus.stopwords}}
\end{frame}

\subsection[Lemmatisation]{Lemmatisation}

\begin{frame}[fragile]{Lemmatisation}
	\begin{itemize}
		\item La \textbf{lemmatisation} est une technique qui consiste à remplacer un mot, par le mot que l'on trouverait dans un dictionnaire. Par exemple un mot pluriel deviendrait singulier. Un verbe conjugué prendrait une forme infinitive...
		\pause
		\item La tokenisation et les mots vides semblaient plus simples, la lemmatisation nécessite d'utiliser des bases de données pour connaitre les différents \textbf{lemmes} et leurs variantes.
		\pause
		\item Une autre difficulté est que la lemmatisation peut dépendre du contexte grammatical. Et donc de la langue ! C'est pour cette raison que l'on se repose sur des librairies qui implémentent tout ça pour nous, comme NLTK.
		\pause
		\item L'analyse grammaticale (\textbf{Part of Speech}) est réalisée par un outil. D'ailleurs il arrive que ce soit lui-même un outil d'apprentissage automatique, comme un réseau de neurones ou une chaine de Markov !
	\end{itemize}
\end{frame}

\begin{frame}[fragile]{Lemmatisation - POS}
	Part of Speech appliqué à notre exemple:
	\begin{lstlisting}
		[('i', 'NN'), ('loved', 'VBD'), ('every', 'DT'), ('second', 'NN'), ('of', 'IN'), ('this', 'DT'), ('masterpiece', 'NN'), ('!', '.')]
		[('the', 'DT'), ('actors', 'NNS'), ('were', 'VBD'), ('terrible', 'JJ'), ('and', 'CC'), ('the', 'DT'), ('story', 'NN'), ('was', 'VBD'), ('boring', 'VBG'), ('.', '.')]
	\end{lstlisting}
	Où nous retrouvons des noms singulier (NN) ou pluriel (NNS), des verbes (VBD), -parfois au gérondif (VBG)-, des déterminants (DT), de la ponctuation (.), des adjectifs (JJ), des conjonctions de coordination (CC). \\
	\vfill
	Méthodes: \href{https://www.nltk.org/api/nltk.stem.wordnet.html#nltk.stem.wordnet.WordNetLemmatizer}{\texttt{nltk.stem.wordnet.WordNetLemmatizer}}
\end{frame}

\begin{frame}[fragile]{Lemmatisation - POS}
	Après une lemmatisation utilisant POS:
	\begin{lstlisting}
		['love', 'second', 'masterpiece']
		['actor', 'terrible', 'story', 'bore']
		['amazing', 'music', 'great', 'special', 'effect']
		['rather', 'watch', 'paint', 'dry', 'see', 'movie']
		['complete', 'waste', 'time', 'money']
		["n't", 'miss', 'masterpiece']
	\end{lstlisting}
	Les mots sont maintenant uniquement singuliers et les verbes ont été remplacés par leur forme infinitive.
\end{frame}

\subsection[Vectorisation]{Vectorisation}

\begin{frame}{Vectorisation - Idée générale}
	Une fois les critiques tokénisées et lemmatisées, il est nécessaire de convertir ces données en \textbf{données numériques} pour qu'elles puissent être traitées par un algorithme d'apprentissage automatique. Cette opération peut être réalisée par une opération simple de \textbf{vectorisation par comptage}:
	\begin{itemize}
		\item L'idée consiste à établir une liste de $d$ tokens qui vont définir un \textbf{vocabulaire}.
		\item Chaque mot dans le vocabulaire est associé à un élément $i$ d'un vecteur.
		\item Chaque critique ayant été convertie en une liste de tokens, on compte combien de fois chaque token du vocabulaire est représenté dans la liste de tokens de la critique.
	\end{itemize}	
\end{frame}

\begin{frame}{Vocabulaire et Corpus}
	\begin{block}{Définition: vocabulaire}
		On appelle \textbf{vocabulaire} l'ensemble $V$ des $d$ tokens (ou mots) \textbf{uniques} présents dans les données d'apprentissage $\mathcal{D}$:
		\begin{equation}
			V = \{m_1, m_2, \cdots, m_d\}
			\nonumber
		\end{equation}
	\end{block}
	\pause
	\begin{block}{Définition: corpus}
		On appelle \textbf{corpus} l'ensemble $\mathcal{D}$ des $n$ documents textuels $t \in \mathcal{T}$  présents dans les données d'apprentissage:
		\begin{equation}
			\mathcal{D} = \{t^{(1)}, t^{(2)}, \cdots, t^{(n)}\}
			\nonumber
		\end{equation}
	\end{block}	
\end{frame}

\begin{frame}{Vectorisation par comptage - Formulation}
	\begin{block}{Vectorisation par comptage}
		La \textbf{vectorisation par comptage} consiste en une fonction $\Phi: \mathcal{T} \to \mathbb{N}^d$ qui permet de déterminer, pour un texte $i$, un vecteur $x^{(i)} \in \mathbb{N}^d$:
		\begin{equation}
			x^{(i)} = \Phi(t^{(i)}) = \begin{pmatrix}
				\text{tf}^{(i)}_1 &	\text{tf}^{(i)}_2 &	\cdots & \text{tf}^{(i)}_d
			\end{pmatrix}
			\nonumber
		\end{equation}
		Où $\text{tf}^{(i)}_j$ représente le nombre de fois (i.e. \textbf{fréquence du mot}) où le mot $m_j$ apparait dans le document textuel $t^{(i)}$. $\text{tf}^{(i)}_j$ peut aussi être décrite avec une \textbf{fonction indicatrice}:
		\begin{equation}
			\text{tf}^{(i)}_j = \underset{w \in t^{(i)}}{\sum} 1(w = m_j)
			\nonumber
		\end{equation}
		La fonction $1(u)$ vaut 1 si $u$ est une condition vraie et 0 sinon.
	\end{block}
	\vfill
	Méthodes: \href{https://scikit-learn.org/stable/modules/generated/sklearn.feature_extraction.text.CountVectorizer.html}{\texttt{sklearn.feature\_extraction.text.CountVectorizer}}
\end{frame}

\begin{frame}{Vectorisation avec TF-IDF}
	Le comptage simple donne trop d'importance aux mots fréquents. Le \textbf{TF-IDF} permet de pénaliser les mots communs à tout le corpus $\mathcal{D}$:
	\begin{block}{Calcul du score TF-IDF}
		Le score pour un mot $m_j$ dans un document $t^{(i)}$ est :
		\begin{equation}
			\text{tf-idf}_j^{(i)} = \text{tf}_j^{(i)} \times \ln\left( \frac{n}{n_{m_j}} \right)
			\nonumber
		\end{equation}
		Avec $n_{m_j}$ le nombre de documents contenant le mot $m_j$.
	\end{block}
	\begin{itemize}
		\item \textbf{TF}: Fréquence du mot dans le document ($\text{tf}^{(i)}_j$).
		\item \textbf{IDF}: Logarithme de l'inverse de la fréquence documentaire. Nul si $n_{m_j} = n$
	\end{itemize}
	\vfill
	Méthodes: \href{https://scikit-learn.org/stable/modules/generated/sklearn.feature_extraction.text.TfidfTransformer.html}{\texttt{sklearn.feature\_extraction.text.TfidfTransformer}} et \href{https://scikit-learn.org/stable/modules/generated/sklearn.feature_extraction.text.TfidfVectorizer.html}{\texttt{sklearn.feature\_extraction.text.TfidfVectorizer}}.
\end{frame}

\subsection[Calcul des probabilités]{Calcul des probabilités}

\begin{frame}{Calcul des probabilités}
	Le vocabulaire $V$ est le \textbf{même} pour toutes les classes. La vectorisation se fait pour chaque classe $k$ séparément à partir des données de la classe $\mathcal{D}_k$. Les probabilités sont calculées par l'algorithme \href{https://scikit-learn.org/stable/modules/generated/sklearn.naive_bayes.MultinomialNB.html}{\texttt{sklearn.naive\_bayes.MultinomialNB}}:
	\begin{block}{Calcul des probabilités avec lissage de Laplace}
		Les probabilités de la loi multinomiale sont calculées séparément pour chaque classe $k$ :
		\begin{equation}
			P(m_j \mid C = k) = \frac{1 + \underset{t^{(i)} \in \mathcal{D}_k}{\sum} \text{tf}^{(i)}_j}{d + \underset{t^{(i)} \in \mathcal{D}_k}{\sum} \sum_{j=1}^d \text{tf}^{(i)}_j}
			\nonumber
		\end{equation}
		Dans le cas où TF-IDF est utilisé, on remplace $\text{tf}^{(i)}_j$ par $\text{tf-idf}^{(i)}_j$.
	\end{block}
	\textit{Note : Le lissage ajoute virtuellement un document de $d$ tokens contenant chaque mot une fois.}
\end{frame}

